\section{Auditoria y trazabilidad}

\subsection{Descripcion}

Este documento define que eventos deben quedar auditados y que informacion minima necesita cada registro para soportar operacion, revision y resolucion de disputas.

\subsection{Alcance}

Aplica a acciones administrativas, cambios de estado relevantes, overrides y procesos automaticos sensibles.

\subsection{Definiciones}

\begin{itemize}
  \item \textbf{Evento auditable}: accion o proceso cuya ocurrencia debe poder demostrarse despues.
  \item \textbf{Actor}: usuario o proceso responsable de un evento.
\end{itemize}

\subsection{Reglas}

Eventos que deben auditarse como minimo:

\begin{itemize}
  \item aprobacion o rechazo de usuario;
  \item bloqueo de usuario;
  \item carga, revision y validacion de comprobante;
  \item cambio de kickoff;
  \item ingreso manual o correccion de resultado;
  \item recalculo de puntajes;
  \item archivado de torneo;
  \item errores operativos relevantes de integracion de resultados.
\end{itemize}

Campos minimos del evento auditable:

\begin{itemize}
  \item identificador del evento;
  \item fecha y hora;
  \item actor;
  \item tipo de accion;
  \item entidad afectada;
  \item valor anterior cuando aplique;
  \item valor nuevo cuando aplique;
  \item motivo o justificacion cuando aplique.
\end{itemize}

\subsection{Casos borde}

\begin{itemize}
  \item Los procesos automaticos deben registrarse con identidad de proceso y no como usuario anonimo.
  \item Si una accion falla despues de un cambio parcial, la auditoria debe reflejar el resultado real de la operacion.
\end{itemize}

\subsection{Invariantes}

\begin{itemize}
  \item Ninguna correccion critica del dominio puede ocurrir sin rastro auditable.
  \item La auditoria debe permitir reconstruir que paso, quien lo hizo y cuando ocurrio.
\end{itemize}

\subsection{Preguntas abiertas}

\begin{itemize}
  \item Confirmar politica de retencion y consulta de eventos de auditoria para uso administrativo.
\end{itemize}

\subsection{Decisiones pendientes}

\begin{itemize}
  \item `Decision pendiente` `No bloqueante`: definir ventana de retencion de eventos auditables y nivel de acceso a consulta historica.
\end{itemize}
